% Generated draft: requires mathematical and source review.
\documentclass[11pt]{article}
\usepackage[margin=25mm]{geometry}
\usepackage{fontspec}
\setmainfont{Latin Modern Roman}
\usepackage{amsmath,amssymb,mathtools,unicode-math}
\setmathfont{Latin Modern Math}
\usepackage{enumitem,needspace,xurl,hyperref}
\setlength{\parindent}{0pt}
\setlength{\parskip}{5pt}
\setlength{\emergencystretch}{3em}
\newcommand{\N}{\mathbb{N}}
\newcommand{\Z}{\mathbb{Z}}
\newcommand{\Q}{\mathbb{Q}}
\newcommand{\R}{\mathbb{R}}
\newcommand{\C}{\mathbb{C}}
\newcommand{\F}{\mathbb{F}}
\newcommand{\A}{\mathbb{A}}
\newcommand{\PP}{\mathbb{P}}
\newcommand{\E}{\mathbb{E}}
\newcommand{\eps}{\varepsilon}
\newcommand{\dd}{\,\mathrm{d}}
\begin{document}
\section*{UnsolvedMath ID: 30002670}
\textbf{Source problem code:} catalogue entry
\subsection{Definitions and mathematical statement}
The supplied record defines a \emph{$k$-gon} to be a convex polygon with at most $k$ vertices. An ordered arrangement is a tuple $(P_1,\ldots,P_n)$, with the labels retained when taking equivalence classes. The source asks whether the space of these arrangements in convex position, modulo affine transformations, is contractible.

A conditional formalization is as follows. If the intended ambient space is $\R^2$ and polygons are compact convex sets with nonempty interior, put
\[
\mathcal P_{\leq k}=\{P\subseteq\R^2:P\text{ is a compact convex polygon with nonempty interior and at most }k\text{ vertices}\},
\]
and let
\[
X_{n,k}=\{(P_1,\ldots,P_n)\in\mathcal P_{\leq k}^{\,n}:(P_1,\ldots,P_n)\text{ is in convex position}\}.
\]
Here ``in convex position'' remains an unresolved source condition: the supplied record does not define it. If all invertible affine transformations of the plane are intended, let
\[
G=\operatorname{Aff}(\R^2)=\{x\mapsto Ax+b:A\in\operatorname{GL}_2(\R),\ b\in\R^2\}
\]
act simultaneously on the polygons, without permuting their labels. One possible topology on $X_{n,k}$ is the subspace topology induced by the product Hausdorff metric on compact sets, where
\[
d_H(P,Q)=\max\left\{\sup_{p\in P}\inf_{q\in Q}\|p-q\|,\ \sup_{q\in Q}\inf_{p\in P}\|q-p\|\right\}.
\]
Under this interpretation, give $Y_{n,k}=X_{n,k}/G$ the quotient topology. The intended question would then be whether, for every admissible pair $(n,k)$, there exist $y_0\in Y_{n,k}$ and a continuous map
\[
H:Y_{n,k}\times[0,1]\longrightarrow Y_{n,k},\qquad H(y,0)=y,\quad H(y,1)=y_0\quad(y\in Y_{n,k}).
\]
This is a conditional formulation template. The supplied record does not specify the definition of convex position, the topology, admissible parameter ranges, or degeneracy conventions. In particular, the planar, nonempty-interior and Hausdorff-topology choices above are explicitly provisional. The usual nondegenerate planar range would be $n\geq1$ and $k\geq3$, but this range is not stated in the record. These qualifications must be resolved before the template can be treated as a precise source-faithful statement.
\subsection{Short English statement}
The question asks whether the space of labelled arrangements of $n$ convex polygons, each with at most $k$ vertices and satisfying the intended convex-position condition, can be continuously contracted to a point after identifying arrangements related by a common affine transformation. The supplied record leaves essential conventions undefined, so the displayed formulation is conditional and requires mathematical review.
\subsection{Sources}
\begin{itemize}
\item Local UnsolvedMath registry, supplied complete record: permanent ID $30002670$, title ``Contractibility of Convex-Position Polygon Spaces,'' source problem code \texttt{catalogue entry}, collection ``Oberwolfach Reports Open Problems.'' No individual proposer is supplied; the proposed year is recorded as $2014$.
\item Supplied source citation: \emph{Discrete Geometry (2015)}, Oberwolfach Reports, DOI 10.4171/owr/2014/40, pp.\ 2235--2296. \url{https://doi.org/10.4171/owr/2014/40}. The background describes Oberwolfach Reports as records of mathematical workshops and their open-problem discussions. The citation's year $2015$ differs from the recorded proposed year $2014$ and the year appearing in the DOI; this discrepancy is retained without resolution.
\item Supplied literature reference: Dobbins--Holmsen--Hubard, ``convex-body realization spaces.'' \url{https://doi.org/10.1007/s00454-017-9869-8}.
\item Supplied literature reference: ``Oberwolfach report.'' \url{https://ems.press/content/serial-article-files/46530}.
\item Literature assessment dated exactly \texttt{2026-08-21}: the supplied curation labels the question open and states that realization spaces of unrestricted convex-body arrangements are understood, that bounded-vertex polygon realization spaces can have complicated topology, and that the convex-position $k$-gon contractibility question remains open. These are supplied dated curation claims, not independently verified conclusions. The background expressly requests independent verification of time-sensitive status and cited claims.
\item The registry labels the statement \texttt{corrected\_verified} and reports that the official report passage was checked directly and its question, definitions and notation restored without strengthening the claim. This is supplied verification metadata; no independent source inspection or mathematical verification is asserted here. The original wording omits ``ordered,'' whereas the cleaned statement and statement field explicitly include it.
\end{itemize}
\end{document}
